\sgaStableTarget{sga6.e10.sec.7.16}\subsection*{7.16. Passage to geometric fibers.}
Let @@M0@@ be an algebraic closure of @@M1@@, the inductive limit of its finite subextensions @@M2@@, and let @@M3@@ be the normalization of @@M4@@ in @@M5@@, the inductive limit of its normalizations @@M6@@ in the @@M7@@. Choose a maximal ideal @@M8@@ of @@M9@@, thus inducing on each @@M10@@ a maximal ideal @@M11@@, and set @@M12@@, which by Krull--Akizuki is again a discrete valuation ring, whose spectrum we denote by @@M13@@ and whose points we denote by @@M14@@ and @@M15@@. The residue field @@M16@@ is an algebraic closure of @@M17@@, and is the inductive limit of the residue fields @@M18@@ of the @@M19@@. Set @@M20@@, @@M21@@. \emph{Now suppose} @@M22@@ \emph{of finite type over} @@M23@@; then @@M24@@ and @@M25@@ are of finite type over @@M26@@ and @@M27@@, respectively, and hence \emph{noetherian}. By IV 3.2.3, we have
@@M28@@
On the other hand, by 7.15 there is a homomorphism from the inductive system of the @@M29@@ to the inductive system of the @@M30@@. Passing to the limit, one obtains a canonical homomorphism:
\sgaStableTarget{sga6.e10.eq.7-16-1}@@M31@@
again called, as expected, the \emph{specialization homomorphism}. It induces a homomorphism @@M32@@ by 7.4, using the fact (whose verification is left to the reader) that the functor @@M33@@ (just like @@M34@@) commutes with the limits occurring here. Similarly, when @@M35@@ is smooth, one can interpret (7.16.1) as a homomorphism on @@M36@@, which is moreover a ring homomorphism, and even a homomorphism of @@M37@@-rings when @@M38@@ is separated; in this case it therefore induces a homomorphism on @@M39@@.

